\documentclass[UTF8,12pt]{ctexart}
\usepackage[a4paper,margin=24mm]{geometry}
\usepackage{amsmath,amssymb,bm,graphicx,adjustbox}
\newcommand{\fitmath}[1]{\begin{adjustbox}{max width=\linewidth}$\displaystyle #1$\end{adjustbox}}
\setlength{\parindent}{0pt}
\setlength{\parskip}{.7em}
\sloppy
\allowdisplaybreaks
\begin{document}
\section*{2011 年考研数学三 · 参考答案}
\subsection*{原 PDF 第 45 页}

\subsection*{2011年真题参考答案}

\subsection*{一、选择题}

（1）C。 （2）B。 （3）A。 （4）B。 （5）D。 （6）C。 （7）D。 （8）D。


\subsection*{二、填空题}

（9）\(\displaystyle (1+3x)e^{3x}\)。 （10）\(\displaystyle (1+2\ln2)(dx-dy)\)。 （11）\(\displaystyle y=-2x\)。 （12）\(\displaystyle \dfrac{\displaystyle 4}{\displaystyle 3}\pi\)。 （13）\(\displaystyle 3y_1^2\)。 （14）\(\displaystyle \mu\sigma^2+\mu^3\)。


\subsection*{三、解答题}

（15）\(\displaystyle -\dfrac{\displaystyle 1}{\displaystyle 2}\)。


（16）\(\displaystyle \left.\dfrac{\displaystyle \partial^2z}{\displaystyle \partial x\partial y}\right|_{(1,\allowbreak 1)}=f^{\prime\prime}_{11}(2,\allowbreak 2)+f^{\prime}_2(2,\allowbreak 2)f^{\prime\prime}_{12}(1,\allowbreak 1)\)。


（17）\(\displaystyle 2\sqrt{x}\arcsin\sqrt{x}+2\sqrt{x}\ln x+2\sqrt{1-x}-4\sqrt{x}+C\)，其中 \(\displaystyle C\) 为任意常数。


（18）证明略。


（19）\(\displaystyle f(x)=\dfrac{\displaystyle 4}{\displaystyle (2-x)^2},\allowbreak \ 0\le x\le1\)。


（20）（I）\(\displaystyle a=5\)。
（II）\(\displaystyle \boldsymbol\beta_1=2\boldsymbol\alpha_1+4\boldsymbol\alpha_2-\boldsymbol\alpha_3,\allowbreak \ \boldsymbol\beta_2=\boldsymbol\alpha_1+2\boldsymbol\alpha_2,\allowbreak \ \boldsymbol\beta_3=5\boldsymbol\alpha_1+10\boldsymbol\alpha_2-2\boldsymbol\alpha_3\)。


（21）（I）特征值 \(\displaystyle -1,\allowbreak 1,\allowbreak 0\)，分别对应特征向量 \(\displaystyle k_1(1,\allowbreak 0,\allowbreak -1)^{\mathrm T},\allowbreak k_2(1,\allowbreak 0,\allowbreak 1)^{\mathrm T},\allowbreak k_3(0,\allowbreak 1,\allowbreak 0)^{\mathrm T}\)，其中 \(\displaystyle k_1,\allowbreak k_2,\allowbreak k_3\) 为任意非零常数。
（II）


\[\fitmath{\displaystyle \begin{pmatrix}0&0&1\\0&0&0\\1&0&0\end{pmatrix}}\]


（22）（I）\(\displaystyle (X,\allowbreak Y)\) 的概率分布为


\[\fitmath{\displaystyle \begin{array}{c|ccc}X\backslash Y&-1&0&1\\\hline 0&0&\dfrac{\displaystyle 1}{\displaystyle 3}&0\\1&\dfrac{\displaystyle 1}{\displaystyle 3}&0&\dfrac{\displaystyle 1}{\displaystyle 3}\end{array}}\]


（II）\(\displaystyle Z=XY\) 的概率分布为


\[\fitmath{\displaystyle \begin{array}{c|ccc}Z&-1&0&1\\\hline P&\dfrac{\displaystyle 1}{\displaystyle 3}&\dfrac{\displaystyle 1}{\displaystyle 3}&\dfrac{\displaystyle 1}{\displaystyle 3}\end{array}}\]


（III）\(\displaystyle 0\)。


（23）（I）


\[\fitmath{\displaystyle f_X(x)=\begin{cases}x,&0<x\le1,\\2-x,&1<x\le2,\\0,&\text{其他}.\end{cases}}\]


（II）


\[\fitmath{\displaystyle f_{X\mid Y}(x\mid y)=\begin{cases}\dfrac{\displaystyle 1}{\displaystyle 2-2y},&y\le x\le2-y,\\0,&\text{其他}.\end{cases}}\]


\clearpage
\end{document}
